\lecture{9}{2026-09-28}{Dominant singularities and \texorpdfstring{\(\mathbb{N}\)}{N}-rational functions}

For a rational generating function, determining coefficient asymptotics is straightforward in two favorable cases:
\begin{itemize}
    \item there is a unique root of the denominator closest to the origin; or
    \item among the roots of minimal modulus, there is a unique root of maximal multiplicity.
\end{itemize}
This raises two natural algorithmic questions: how do we identify the roots of minimal modulus, and what is the computational complexity of extracting dominant asymptotics?


\section*{The Vivanti--Pringsheim theorem}

When a sequence has non-negative coefficients, the search for dominant singularities is greatly simplified by locating real positive singularities.

\begin{theorem}[Vivanti--Pringsheim Theorem, {\cite[Proposition 2.4]{Melczer2021}}]
    \label{thm:vivanti-pringsheim}
    Let \(F(z) = \sum_{n \ge 0} f_n z^n\) have radius of convergence \(R > 0\).
    If \(f_n \ge 0\) for all \(n \ge 0\), then \(z = R\) is a singularity of \(F(z)\).
    In other words, at least one singularity of minimal modulus is a positive real number.
\end{theorem}

\begin{proof}
    Since \(f_n \ge 0\) for all \(n\), the triangle inequality yields
    \[
        |F(z)| = \left| \sum_{n=0}^{\infty} f_n z^n \right| \le \sum_{n=0}^{\infty} f_n |z|^n = F(|z|)
    \]
    for all \(z\) in the open disk of convergence \(\{|z| < R\}\).
    If \(\omega\) is a polar singularity of \(F(z)\) of minimal modulus \(|\omega| = R\), then approaching \(\omega\) from within the disk of convergence yields \(|F(z_k)| \to \infty\) along some sequence \(z_k \to \omega\).
    Because \(F(|z_k|) \ge |F(z_k)|\), it follows that \(F(|z_k|) \to \infty\) as \(|z_k| \to R\).
    Thus, \(z = R = |\omega|\) must be a singularity of \(F(z)\).
\end{proof}


\section*{Decidability problems for \(C\)-finite sequences}

Although rational functions admit explicit formulas for their coefficients, basic algorithmic decision problems about \(C\)-finite sequences remain surprisingly deep.

\begin{remark}[Skolem's problem, {\cite[Open Problem 2.1]{Melczer2021}}]
    Is there an algorithm that takes a \(C\)-finite sequence \((f_n)\) over \(\mathbb{Q}\) and decides whether there exists an index \(n\) such that \(f_n = 0\)?
    \begin{itemize}
        \item Recurrences of order \(1\) and \(2\) are trivial.
        \item Recurrences of orders \(3\) and \(4\) were proven decidable in the 1980s by Mignotte et al.\ and Vereshchagin.
        \item Recurrences of order \(5\) and higher remain open in general.
        For simple recurrences, decidability for prime power indices was established in 2022 (Kenison et al.) under standard number-theoretic conjectures.
    \end{itemize}
\end{remark}

\begin{remark}[Ultimate positivity, {\cite[Open Problem 2.2]{Melczer2021}}]
    Does there exist an algorithm that takes a \(C\)-finite sequence \((f_n)\) over \(\mathbb{Q}\) and decides whether \(f_n \ge 0\) for all but finitely many \(n\) (i.e., whether \((f_n)\) is eventually non-negative)?
    \begin{itemize}
        \item Decidability is known for rational functions with denominators of degree at most \(5\), and for arbitrary degree when the denominator is square-free (Ouaknine and Worrell).
        \item Degree \(6\) and above remains open in general.
    \end{itemize}
\end{remark}


\section*{\texorpdfstring{$\mathbb{N}$}{N}-rational functions}

Combinatorial sequences naturally avoid such pathological undecidability behavior because their generating functions belong to a well-behaved subclass of rational series.

\begin{definition}[\(\mathbb{N}\)-rational functions, {\cite[Definition 2.8]{Melczer2021}}]
    \label{def:n-rational}
    The set of \emph{\(\mathbb{N}\)-rational functions} is the smallest set of functions:
    \begin{itemize}
        \item containing the constant \(1\) and the variable \(z\);
        \item closed under addition and multiplication;
        \item closed under the quasi-inverse operation \(f \mapsto \frac{1}{1 - zf}\).
    \end{itemize}
\end{definition}

In formal language theory, \(\mathbb{N}\)-rational functions are precisely the generating functions of regular languages counted by word length.

\begin{theorem}[Berstel's Theorem, {\cite[Proposition 2.6]{Melczer2021}}]
    \label{thm:berstel}
    If \(F(z)\) is an \(\mathbb{N}\)-rational function, then every dominant singularity of \(F(z)\) has the form \(\rho \zeta\), where \(\rho > 0\) is the radius of convergence and \(\zeta\) is a root of unity.
    In particular, all dominant singularities are contained in a set \(\{\rho, \rho \zeta_1, \dots, \rho \zeta_s\}\) with \(\zeta_j\) roots of unity.
\end{theorem}

\begin{example}[A non-\(\mathbb{N}\)-rational series with non-negative coefficients, {\cite[Example 2.6]{Melczer2021}}]
    Let \(\theta = \arccos(1/3)\), and consider the rational function
    \begin{align*}
        F(z) &= \frac{2(1+3z)^2}{(1-9z)(1+14z+81z^2)} \\
        &= \frac{1/2}{1 - 9z e^{2i\theta}} + \frac{1/2}{1 - 9z e^{-2i\theta}} + \frac{1}{1 - 9z}.
    \end{align*}
    The dominant singularities are \(z = 1/9\) and \(z = \frac{1}{9} e^{\pm 2i\theta}\).
    Because \(e^{2i\theta}\) is not a root of unity, \(F(z)\) is not \(\mathbb{N}\)-rational by \Cref{thm:berstel}.
    Nevertheless, its power series coefficients satisfy
    \[
        [z^n] F(z) = 9^n \left( 1 + \cos(2n\theta) \right) \ge 0,
    \]
    and are integers because the denominator of \(F(z)\) has integer coefficients with constant term \(1\).
\end{example}

\begin{corollary}[Multisection of \(\mathbb{N}\)-rational functions]
    \label{cor:n-rational-multisection}
    If \(F(z)\) is \(\mathbb{N}\)-rational, then there exists an integer \(p \ge 1\) such that \(F(z)\) can be decomposed as
    \[
        F(z) = S_0(z^p) + z S_1(z^p) + \dots + z^{p-1} S_{p-1}(z^p),
    \]
    where each \(S_0, \dots, S_{p-1}\) is \(\mathbb{N}\)-rational and has a unique dominant singularity.
\end{corollary}

\begin{remark}[Combinatorial meta-theorem]
    Every ``naturally occurring'' rational generating function arising from a combinatorial class is \(\mathbb{N}\)-rational.
    Consequently, dominant singularities in combinatorial enumeration differ by roots of unity, and periodic asymptotic behavior can be effectively decoupled via multisection.
\end{remark}


\section*{Root separation and computational complexity}

How efficiently can we extract asymptotics for rational generating functions?
The most computationally intensive step is separating and identifying the roots of the denominator of minimal modulus.

Let \(P(z) \in \mathbb{Z}[z]\) be a polynomial of degree \(d\) whose coefficients have bit size at most \(h\) (height bounded by \(2^h\)).

\begin{lemma}[Root separation bound]
    \label{lem:root-separation}
    If \(\alpha \neq \beta\) are distinct roots of \(P(z)\), then
    \[
        |\alpha - \beta| \ge D_d \, 2^{-h(d-1)},
    \]
    where \(D_d\) is a constant depending on \(d\) that satisfies \(D_d = d^{-O(d)}\).
\end{lemma}

Therefore, distinguishing distinct roots of \(P(z)\) requires on the order of \(O(hd)\) bits of precision.

For coefficient asymptotics, however, we must separate not just distinct roots, but their \emph{moduli}:
\[
    \big| |\alpha| - |\beta| \big| \ge \Delta(d, h) \quad \text{whenever } |\alpha| \neq |\beta|.
\]
Effective lower bounds on root modulus separation are significantly more delicate, and this problem remains an active area of investigation.